% ============================================================
% §1 Introduction — OR journal style
% Expanded motivation, OR framing, contributions as full paragraphs
% Related work merged as §1.1 subsection
% ============================================================
\section{Introduction}
\label{sec:intro}

% --- OR-motivated opening ---
Diffusion models generate data by following a \emph{denoising trajectory}: starting
from random noise, the model applies a sequence of learned noise-reduction steps
until a sample from the target distribution emerges.
Originally developed for image synthesis and now powering large-scale
text-to-image systems such as Stable Diffusion, these models rest on stochastic
process theory, where a forward process progressively corrupts data with noise
and a reverse process learns to undo the corruption step by step.
Their mathematical foundations connect directly to applied probability and
stochastic simulation; see
\citet{song2019generative}, \citet{ho2020denoising}, the comparative survey
by \citet{cao2024comparative}, and the rigorous SDE-based treatment
by \citet{tang2024sde}.
As simulation tools, diffusion models combine stable training, strong mode
coverage, and flexibility for both unconditional and conditional
generation, properties especially valuable in operations research and
management science, where realistic scenario generation from learned
distributions supports downstream optimization, risk assessment, and policy
evaluation.

% --- Problem: inference-time compute allocation ---
While most progress in diffusion models has targeted training objectives and
model architectures, an emerging line of work demonstrates that
\emph{inference-time scaling}, allocating additional computation at test time
without retraining, sharply improves sample quality under a fixed model.
Test-time techniques such as noise search, trajectory refinement, and
verifier-guided sampling yield substantial gains over plain diffusion sampling
when additional compute is available~\citep{ma2024inference,ramesh2025nts}.

Every diffusion-family model, whether a DDPM, a score-based SDE, or a
flow-matching ODE, generates samples by traversing a discrete trajectory of
$T$ timesteps.
Inference-time compute enters this trajectory in two fundamentally different
ways:
(i)~\emph{lengthening the trajectory} by increasing $T$ (more denoising steps),
which yields quickly diminishing returns, and
(ii)~\emph{searching within each timestep} by evaluating multiple noise
candidates and selecting the one that maximizes a quality metric.
The second axis is the focus of this paper: under a fixed budget of function
evaluations (NFE), how should per-timestep search effort be allocated across the
trajectory?

% --- Gap: uniform allocation ---
Existing inference-time methods waste computation by distributing search budget
uniformly across denoising steps, ignoring that different steps differ
dramatically in their sensitivity to noise perturbations.
Most approaches either (i)~allocate the same number of search iterations to
every timestep, or (ii)~hard-code search to a fixed window of steps
(e.g., only the initial noise or a predetermined middle segment).
Refining noise at high-sensitivity steps yields large downstream quality
improvements, while the same computation at low-sensitivity steps has negligible
effect.

% --- Our approach ---
In this work, we study \emph{where to search} along the diffusion trajectory.
We formalize inference-time noise refinement as a \emph{noise trajectory search}
problem under a fixed number of function evaluations, and ask how to
\emph{globally schedule} search computation across timesteps.
The key insight is that effective inference-time scaling requires reasoning at
two distinct levels: \emph{local} decisions about how to search for better noise
at a given timestep, and \emph{global} decisions about how to distribute
computation across the entire denoising process.
While prior work has explored increasingly sophisticated local search operators,
global per-step compute allocation has only recently received attention, with
concurrent work proposing either uniform allocation with implicit
adaptation~\citep{ramesh2025nts} or purely reactive online
reallocation~\citep{kim2025rbf}.

To address this gap, we propose a unified two-level framework for noise
trajectory search in diffusion models.
At the local level, a broad family of existing noise refinement methods reduces
to a generic local search operator applied at a single timestep.
At the global level, a scheduler decides \emph{when} and \emph{how much} local
search to perform at each timestep under a total compute budget.
Within this framework, we show that standard diffusion sampling, best-of-$N$
selection, $\epsilon$-greedy noise search, and particle- or tree-based methods
are special cases with fixed or heuristic scheduling policies.

Building on this formulation, we propose \textbf{GAINS} (\textbf{G}lobal
\textbf{A}daptive \textbf{I}nference-time \textbf{N}oise \textbf{S}cheduling),
a global scheduling algorithm that combines \emph{offline profiling} with
\emph{online control}.
Offline, GAINS profiles verifier sensitivity across timesteps to obtain a coarse
compute allocation tailored to the model and dataset.
Online, GAINS adjusts per-step search using instance-level feedback (score
gains and variance), stopping early on unproductive steps and reallocating
computation to productive ones.
The controller operates causally (no access to future timesteps) and respects a
strict NFE budget.

% --- Contributions (OR style: numbered, full paragraphs) ---
We make the following contributions.

1.~We introduce a unified \emph{two-level view} of noise trajectory search that
separates local noise refinement from global timestep scheduling.
This abstraction covers both stochastic (DDPM/SDE) and deterministic (flow/ODE)
samplers, providing a common language in which existing inference-time search
methods reduce to special cases of scheduling policies (\cref{sec:framework}).

2.~We develop GAINS, a \emph{global scheduling algorithm} that integrates
offline timestep profiling with online feedback control.
Unlike purely uniform~\citep{ramesh2025nts}, purely
reactive~\citep{kim2025rbf}, or fixed-phase~\citep{zhang2025tscend}
approaches, GAINS combines learned per-step difficulty patterns with adaptive
per-instance early stopping, enforcing an exact NFE budget
(\cref{sec:algorithm}).

3.~We establish \emph{theoretical foundations} grounded in the diffusion SDE.
A second-order Taylor expansion of the verifier score function shows that
the per-step score distribution is approximately location-scale with
sensitivity $\sigma_t = g_t\sqrt{h}\,\|\nabla\Psi_t\|$, the product of the
diffusion coefficient, step size, and verifier gradient norm
(\cref{thm:loc-scale}).
This explicit formula explains why timestep sensitivity varies and yields a
principled allocation: optimal offline budgeting follows a water-filling
structure (\cref{prop:offline}), and any fixed allocation incurs a Jensen gap
that online adaptation can partially recover (\cref{prop:online}).

4.~We evaluate GAINS on both text-to-image (Stable Diffusion) and
class-conditional (EDM) generation, as well as flow-based models,
showing consistent improvements over existing inference-time methods under
identical NFE budgets.
GAINS matches or exceeds the uniform-allocation benchmark with 20--50\% fewer function
evaluations; the gains hold across local search operators, prompt sets, and
model families (\cref{sec:experiments}).

% ============================================================
% §1.1 Related Literature (OR convention: subsection of Introduction)
% ============================================================
\subsection{Related Literature}
\label{sec:related-lit}

\paragraph{Inference-time scaling for diffusion models.}
Inference-time scaling improves sample quality by spending additional
computation at test time without retraining~\citep{ma2024inference,ramesh2025nts}.
\citet{ma2024inference} systematize this direction by proposing a framework
organized around two design axes, the \emph{verifier} (any external quality
metric) and the \emph{search algorithm} (best-of-$N$, zero-order search, or
path search), but apply each algorithm with a fixed, per-timestep-uniform
strategy and do not learn which timesteps benefit most from search.
GAINS extends this line by focusing on the complementary question of
\emph{how to allocate} search compute across the trajectory.

\paragraph{Noise search and trajectory refinement.}
Several methods perform local search over noise variables during diffusion
sampling.
Noise quality varies markedly across timesteps and samples:
\citet{qi2024noises} show that not all initial noises suit generation equally
and propose noise selection via inversion stability, while
\citet{zhou2025golden} learn a noise prompt network that transforms random
Gaussian noise into ``golden noise'' optimized for text--image alignment.
These results motivate per-step noise refinement at inference time.
Best-of-$N$ methods sample multiple full trajectories and select the best final
output, but do not refine noise within a trajectory~\citep{ma2024inference}.
\citet{tang2024dno} formulate noise refinement as continuous optimization
(Direct Noise Optimization), applying gradient- or zero-order updates at every
timestep, again with uniform per-step effort.
\citet{ramesh2025nts} cast diffusion as an MDP and relax it to a sequence of
independent contextual bandits, proposing an $\epsilon$-greedy algorithm that
implicitly adapts exploration vs.\ exploitation across timesteps; they note
in an appendix that varying $K_t$ across timesteps can reduce NFE by half
with minimal quality loss, but leave the design of an adaptive $K_t$ scheduler
as future work.
\citet{kim2025rbf} introduce Rollover Budget Forcing (RBF), which rolls over
unused per-step budget to subsequent timesteps, a purely online, reactive
mechanism without prior knowledge of timestep importance.
RBF is developed and validated exclusively on flow models (FLUX),
relying on ODE-to-SDE conversion for particle sampling; its applicability
to DDPM or SDE-based diffusion models remains unvalidated.
\citet{zhang2025tscend} propose T-SCEND, a two-phase strategy that
applies best-of-$N$ in early steps and MCTS in late steps; however, the
phase split is fixed and hand-designed rather than data-driven.
More recently, \citet{zhang2025classical} combine local Langevin MCMC with
global tree search (BFS/DFS) for inference-time control, and
\citet{lee2025abcd} introduce Adaptive Bi-directional Cyclic Diffusion (ABCD),
which dynamically adjusts exploration depth per instance.
\citet{ttsnap2025} propose TTSnap, which prunes low-quality noise candidates
early without full denoising, reallocating saved compute to explore more
candidates; however, this operates at the \emph{candidate} level within a
single timestep rather than scheduling budget \emph{across} timesteps.
In the particle-based setting, \citet{singhal2025fk} introduce a Feynman-Kac
steering framework that allocates non-uniform particle budgets across denoising
stages, and \citet{kim2025psi} propose $\Psi$-Sampler for SMC-based
inference-time alignment with improved initial particle sampling.
\citet{eyring2025hypernetworks} take an orthogonal approach by amortizing
noise search into a learned Noise Hypernetwork, replacing test-time
optimization with a single forward pass.
While these concurrent methods introduce various forms of non-uniform
allocation, they are either purely reactive (RBF, ABCD), use fixed phase
splits (T-SCEND), operate at the candidate rather than timestep level (TTSnap),
or use SMC resampling rather than search-based allocation (FK Steering,
$\Psi$-Sampler).
GAINS explicitly combines \emph{offline profiling}, learning which
timesteps are inherently more amenable to search, with \emph{online
control} that adapts per-step effort to each prompt, achieving non-uniform
budget allocation with strict NFE guarantees.

\paragraph{Adaptive compute allocation.}
Adaptive computation has been studied extensively in other domains, such as early
exiting in neural networks, adaptive depth inference, and budgeted optimization.
In diffusion models, \citet{zhang2025adadiff} propose AdaDiff for adaptive
denoising step selection conditioned on input complexity,
\citet{ye2025tpdm} learn a per-instance time prediction module via RL that
decides which timesteps to visit (adapting the \emph{schedule} rather than the
per-step \emph{search budget}), and
\citet{yeh2025demons} introduce Sampling Demons, a training-free alignment
method that reweights per-step noise importance to steer generation toward
high-reward regions.
These techniques, however, primarily modify the base sampler or noise
distribution rather than allocating additional \emph{search} computation
across timesteps.
\citet{sundaresha2025vt} propose the Verifier-Threshold method for more
efficient non-uniform NFE allocation, demonstrating 2--4$\times$ speedups
over RBF; however, their approach is also purely online and does not use
offline profiling.
GAINS is complementary to sampler design: it keeps the underlying sampler
fixed and schedules inference-time search around it, combining a learned offline
allocation with online gain-and-variance feedback.

\paragraph{Resource allocation and simulation optimization.}
Allocating a fixed computational budget across sequential stages to maximize
aggregate performance is a classical discrete resource allocation
problem~\citep{ibaraki1988resource}.
In our setting, the ``stages'' are denoising timesteps and the ``resource'' is
the number of function evaluations available for noise search.
The local search operator, which evaluates $K$ noise candidates and selects the
best, is closely related to \emph{ranking and selection} in simulation
optimization, where a sampling budget is distributed across alternatives to
identify the best~\citep{chen2000simulation}.
\citet{lam2025importance} study a structurally similar problem in stochastic
optimization, jointly optimizing decision variables and importance sampling
distributions via a single-loop algorithm that achieves oracle-matching
variance despite circular dependence between the two.

\paragraph{Theoretical foundations of diffusion models.}
Several recent theoretical advances inform the design of GAINS.
\citet{gao2024reward} develop a continuous-time RL framework for
reward-directed score-based diffusion, showing that the optimal exploration
policy is Gaussian with time-varying variance, a structure consistent with the
location-scale assumption underlying our analysis.
\citet{tang2024contractive} characterize a per-timestep error sensitivity
function $u(t)$ governing how score estimation errors propagate through the
backward process; this function provides a theoretical proxy for the per-step
variance profile $\{\sigma_t\}$ that drives our water-filling allocation.
\citet{zhou2024jumpdiffusion} extend continuous-time RL to jump-diffusion
processes with martingale-based q-learning, a framework that connects to our
MDP formulation in \cref{sec:mdp}.
